\section{まとめ}
\label{sec:conclusion}

本研究では，数値ラベル $\boldsymbol{c}$ を指定するだけでベース動作データから所望の変位だけずれた動作をロボットに実行させるフレームワークを提案した．
このフレームワークでは，ILC による数値収束の前提条件としてNNの単調増加性が必要であり，
``軌道全体の絶対値を予測するよりも残差 $\Delta\xi$ を予測する方がラベルとの単調増加性が強まる'' という仮説を立てた．

この仮説を検証するため，積み木のピックアンドプレースタスクを用いた予備実験を行った．
残差予測モデル \texttt{res} は Spearman 順位相関 $\rho \approx 0.94$，決定係数 $R^2 \approx 0.87$ を達成し，
直接予測モデル（$\rho \approx 0.65$，$R^2 \approx 0.50$）と比べて大幅に高い単調増加性を示した．
これにより，\textbf{仮説を実験的に支持する結果が得られた}．

一方，以下の 2 点が今後の改善課題として明確になった：

\noindent\textbf{(a) $\boldsymbol{c} = [0,0]$ 時の系統バイアス}：
残差モデルでは補正なし状態でも 15〜19\,mm のオフセットが生じており，``残差ゼロ'' を正確に学習できていない．
損失関数への $\boldsymbol{c} = \boldsymbol{0}$ 時のペナルティ追加や，ILC によるバイアスの反復補正が有効と考えられる．

\noindent\textbf{(b) 姿勢干渉}：
\texttt{res} ではピック時の手先姿勢偏差がラベルと相関しており，残差予測が位置変位だけでなく姿勢変化も学習してしまっている．
姿勢成分の残差をゼロに拘束する構造的制約の導入が有効と考えられる．

今後の方針として，双方向Transformerを用いた一括補正モデルへの移行，Broyden 法を用いた ILC ループの実装，
および位置だけでなく力情報を含むタスク（コネクター挿入・ねじ締め等）への展開を予定している．
